\section{Introduction}\label{sec:intro}
\begin{theorem}\label{thm:main}
    There exists an algorithm that, given the classical description of an $n$-qubit pure state $\ket{\psi_{target}}$ and a single copy of an $n$-qubit (possibly mixed) state $\rho_{lab}$, performs only single-qubit measurements on $\rho_{lab}$ and outputs:
    \begin{itemize}
        \item \textsc{Accept} with probability at least $\mathrm{Fid} \coloneqq \bra{\psi_{target}}\rho_{lab}\ket{\psi_{target}}$;
        \item \textsc{Reject} with probability at least $\mathrm{Infid}/n \coloneqq (1-\mathrm{Fid})/n$.
    \end{itemize}
\end{theorem}

\subsection{Tomography (density estimation)}
Before discussing state certification, we start by reviewing the best known results for the baseline problem of \emph{tomography}; i.e., producing an estimate $\wh{\rho}$ of an unknown density matrix $\rho \in \C^{d \times d}$, given $n$ copies $\rho^{\otimes n}$, up to error~$\eps$ (whp) for some notion of ``distance''. We will also let $k$ denote the \emph{rank} of $\rho$, which is $1$ when $\rho$ is a pure state, and in general is at most~$d$.  The best results achievable depend on the ``figure of merit'' --- i.e., distance measure --- chosen (see \Cref{sec:distances} for a review).

In~\cite{HHJ+16} it was shown that  $n = O(kd/\epsilon) \cdot \log(d/\epsilon)$ copies suffice to obtain infidelity~$\eps$ (i.e., $\Fid{\rho}{\wh{\rho}} \leq 1 - \eps$); this also implies that $n = O(kd/\epsilon^2) \cdot \log(d/\epsilon)$ copies suffice to obtain trace distance~$\eps$ (i.e., $\Dtr{\rho}{\wh{\rho}} \leq \eps$).  Those authors also showed that $n = \Omega(kd/\epsilon^2) / \log(d/k\epsilon)$ copies are necessary, with the $\log$ factor being removable in the case $k = d$.  Independently, in~\cite{OW16} it was shown that $n = O(d/\eps^2)$ copies suffice to obtain Hilbert--Schmidt distance~$\eps$ (i.e., $\DHS{\rho}{\wh{\rho}} \leq \eps$); this also implies a copy complexity of $n = O(kd/\epsilon^2)$ for trace distance (slightly better than in~\cite{HHJ+16}).  More generally, \cite{OW16} showed a kind of ``PCA'' result: for $\rho$ of any rank, $n = O(kd/\epsilon^2)$ copies suffice to produce an estimate $\wh{\rho}$ whose trace distance from~$\rho$ is at most $\eps$ more than that of the best rank-$k$ approximator.  Finally, a followup work~\cite{OW17} gave an alternate proof of the $n = O(kd/\epsilon) \cdot \log(d/\epsilon)$ bound for infidelity, showed also an $n = O(k^2d/\epsilon)$ bound, and extended these bounds to the PCA case.

\subsection{Density testing}

Tomography results suffer from the inherent issue that $n = \wt{\Theta}(d^2)$ copies are needed in the general case (except when the figure of merit is Hilbert--Schmidt distance, but this metric is not considered to be very meaningful, operationally).  Thus as mentioned, it is natural to focus on restricted problems like state certification, distance estimation, and other \emph{property testing} problems that can potentially be carried out with $n = O(d)$ or better.  Montanaro and de~Wolf~\cite{MdW16} have given an excellent survey on property testing of quantum states; we review a few of the known results here.

A typical quantum property testing problem would involve two disjoint classes $\calC_1$, $\calC_2$ of $d$-dimensional quantum states; given $n$ copies of an unknown $\rho$, promised to be in either~$\calC_1$ or~$\calC_2$, the task is to distinguish which is the case (whp) using few copies of~$\rho$.  In particular, the \emph{quantum state certification} problem for fixed state $\sigma \in \C^{d \times d}$ is the case when $\calC_1 = \{\sigma\}$ and $\calC_2 = \{ \rho : \mathrm{D}(\rho,\sigma) > \eps\}$ for some notion $\mathrm{D}(\cdot, \cdot)$ of distance and some parameter~$\eps$.

When $\sigma$ is a \emph{pure} state, it is straightforward to show (see, e.g.,~\cite{MdW16}) that the associated quantum state certification task, with infidelity as the distance measure, can be done using $n = O(1/\eps)$ copies (and this implies $n = O(1/\eps^2)$ copies suffice for trace distance).  Indeed, the same is possible when both $\rho$ \emph{and} $\sigma$ are unknown pure states, and one is given $n$ copies of each.  For practical purposes, it may be useful to have a state certification algorithm for a known pure~$\sigma$ that only uses simple measurements; e.g., Pauli observables.  For this problem, it has been shown~\cite{FL11,SLP11,AGKE15} that for $\sigma$ known and pure, one can solve the certification problem given $n = O(d/\eps^2)$ copies of an unknown~$\rho$ with infidelity as the distance metric --- indeed, with this many copies one can estimate the fidelity $\Fid{\rho}{\sigma}$ to $\pm \eps$.

For the state certification problem when $\sigma$ is mixed (not pure), not much is known except in one case: when $\sigma = \Id/d$, the ``maximally mixed'' state.  For this problem, it was shown in~\cite{OW15} that $n = \Theta(d/\eps^2)$ copies are necessary and sufficient, when the distance measure is trace distance.  In fact, for the $n = O(d/\eps^2)$ upper bound, \cite{OW15}~effectively show that one can estimate the \emph{purity} $\tr(\rho^2)$ of~$\rho$ sufficiently well so as to distinguish between purity $1/d$ (achieved by the maximally mixed state) and purity exceeding $1/d + \eps^2/d$.  Note that the latter case is equivalent to $\rho$ being $\epsilon/\sqrt{d}$-far from $\Id/d$ in Hilbert--Schmidt distance and $\epsilon^2$-far from $\Id/d$ in Bures $\chi^2$-divergence. The lower bound was mentioned earlier as \Cref{thm:paninski-lower}.

\subsection{Certification with restricted measurements} \label{sec:restricted-measurement}

The sample-optimal protocol mentioned above implicitly assumes one can implement either the two-outcome POVM $\{\ket{\psi_{target}}\bra{\psi_{target}}, \Id - \ket{\psi_{target}}\bra{\psi_{target}}\}$ or a SWAP test~\cite{BCWdW01} between $\rho_{lab}$ and a freshly prepared copy of $\ket{\psi_{target}}$.  Both options sit awkwardly with the goal of certification: each requires the very capability whose successful realization certification was meant to verify --- namely, preparing the (possibly highly entangled) $n$-qubit pure state $\ket{\psi_{target}}$.

The framework of \emph{classical shadows}~\cite{HKP20} partially side-steps this circularity by replacing target-dependent measurements with random Clifford measurements on independent copies of $\rho$, after which the fidelity $\Fid{\rho}{\ket{\psi_{target}}\bra{\psi_{target}}}$ can be estimated to additive error~$\eps$ using $O(1/\eps^2)$ copies, independent of dimension.  Implementing random Clifford measurements still requires deep circuits that may be impractical, however, and the efficiency of the protocol depends on computing overlaps between $\ket{\psi_{target}}$ and random stabilizer states --- a step that is itself intractable when $\ket{\psi_{target}}$ has high stabilizer rank.

If one further restricts to single-qubit Pauli measurements, the prior protocols of most relevance are the direct fidelity estimation method~\cite{FL11,SLP11,AGKE15} discussed in the previous subsection, and a Pauli-based variant of classical shadows~\cite{HKP20}.  Both have sample complexity that grows exponentially in the worst case --- for direct fidelity estimation, in the system dimension $d = 2^n$; for the Pauli-shadow variant, in the weight of the observable being estimated.  The latter is therefore well suited to certifying small subsystems but cannot efficiently estimate fidelity with a highly entangled $n$-qubit target.

\subsection{Certification of structured state classes and mixed targets} \label{sec:structured-classes}

For target states drawn from structured classes, certification using simple measurements can be substantially cheaper than for arbitrary targets.  Stabilizer states and outputs of shallow quantum circuits admit efficient certification by direct fidelity estimation~\cite{FL11} or by randomized Pauli shadows~\cite{HKP20,SubsystemShadows}.  Efficient protocols are also known for hypergraph states~\cite{HypergraphCert}, output states of IQP circuits~\cite{IQPCert}, bosonic Gaussian states~\cite{BosonicGaussianCert}, and fermionic Gaussian states~\cite{FermionicGaussianCert}.  For generic random target states, cross-entropy benchmarking (XEB) provides a good fidelity estimate under suitable noise models~\cite{XEB1,XEB2,XEB3,XEB4}.

A complementary line of work, extending beyond pure targets, has explored the certification of mixed states~\cite{MixedCert1,MixedCert2,MixedCert3}.

\subsection{The asymptotic regime for state discrimination}  \label{sec:asympt-state-discrim}

There is a related class of work that we refer to as the ``asymptotic regime''.  Consider the simplest quantum property testing problem, \emph{state discrimination}, in which $\calC_1 = \{\sigma_1\}$ and $\calC_2 = \{\sigma_2\}$ for two known states $\sigma_1, \sigma_2 \in \C^{d \times d}$.  The perspective we take in this paper involves determining the least number of copies~$n$ such that one can distinguish $\rho = \sigma_1$ from $\rho = \sigma_2$ with high probability --- say, with both ``type I'' and ``type II'' errors having probability at most~$\delta = 1/3$.  One can reduce this~$\delta$ to any small positive constant at the expense of making~$n$ a constant factor larger.  We refer to this perspective as the \emph{non-asymptotic regime}, because we do not consider any limiting error rate as $n \to \infty$; rather, we  wish to find a concrete upper bound on the~$n$ that suffices, depending only on $d$, the distance between $\sigma_1$ and $\sigma_2$, and nothing else.

On the other hand, there is substantial work on the \emph{asymptotic regime}, sometimes going under the name \emph{quantum hypothesis testing}, in which the focus is on how exponentially fast the error rate goes to~$0$ in the limit as $n \to \infty$.  Here one might seek the best (smallest) constant~$R$ such that, given $n$ copies, one can ensure type~I and type~II errors have probability at most $(R+o(1))^n$, where the $o(1)$ refers to $n  \to \infty$.  A downside of such results is that they do not a priori give any information about how large $n$ needs to be before error bounds ``kick in''; e.g., the $o(1)$ function might not be less than, say,~$.1$ until $n$ is larger than some uncontrolled function of~$d$ (e.g., $2^d$) or of some other parameters (e.g., the smallest nonzero eigenvalue of~$\sigma_1$ or $\sigma_2$).

A good survey of the results in the asymptotic regime appears in~\cite{ANSV08}; they review known quantum versions of Stein's Lemma and Sanov's Theorem, and  prove quantum versions of Chernoff's Bound and the the Hoeffding--Blahut--Csisz\'{a}r--Longo bound.  For example, in the basic hypothesis testing problem described above, they prove that the best rate~$R$ is given by $Q_{\text{min}}\diverg{\sigma_1}{\sigma_2} = \min_{0 \leq s \leq 1}\tr(\sigma_1^s \sigma_2^{1-s})$ (a quantity that is within a factor of~$2$ of the infidelity between $\sigma_1$ and~$\sigma_2$).
